\documentclass[10pt]{article}
\usepackage[margin=1in]{geometry}
\usepackage{amsmath,amssymb}
\usepackage{graphicx}
\usepackage{booktabs}

\title{Appendix: Methods and Proofs}
\author{Mark Rowe Traver}

\begin{document}
\maketitle

\section{Emergence Coordinate Derivation}
\label{sec:appendix_coords}

The four emergence coordinates $(C, F, A, R)$ are linear composites of z-scored features. Let $\mathbf{x} \in \mathbb{R}^{17}$ be the raw feature vector and $\mathbf{z} = (\mathbf{x} - \boldsymbol{\mu}) \oslash \boldsymbol{\sigma}$ be the z-scored vector, where $\boldsymbol{\mu}$ and $\boldsymbol{\sigma}$ are the empirical mean and standard deviation across all systems.

\subsection{Curvature coordinate $C$}

\[
C = \frac{1}{4} \left( -z_{\text{temporal\_corr}} - z_{\text{eff\_rank}} + z_{\text{pc1\_ratio}} - z_{\text{replay}} \right)
\]

$C$ captures the balance between temporal coherence (negative loading on temporal correlation and effective rank) and structural complexity (positive loading on PC1 ratio, negative on replay displacement).

\subsection{Flow coordinate $F$}

\[
F = \frac{1}{4} \left( z_{\tau_{m2}} + z_{\tau_{m4}} + z_{\phi_{\text{corr}}} + z_{pc2} \right)
\]

$F$ captures the amplitude and phase structure of the dynamical flow through the second and fourth moment timescales and the phase-space correlation.

\subsection{Amplitude coordinate $A$}

\[
A = \frac{1}{3} \left( \|\mathbf{z}_{\text{abl}}\|_2 + \sigma(\mathbf{z}_{\text{abl}}) + z_{\tau_{m1}} \right)
\]

where $\mathbf{z}_{\text{abl}} = (z_{\text{abl\_full}}, z_{\text{abl\_no\_m1}}, z_{\text{abl\_no\_m2}}, z_{\text{abl\_no\_m3}}, z_{\text{abl\_no\_m4}}, z_{m2\_\text{contrib}})$.

$A$ captures the amplitude structure through the norm and spread of ablation sensitivities and the first moment timescale.

\subsection{Recurrence coordinate $R$}

\[
R = \frac{1}{3} \left( z_{pc1} + \bar{z}_{\tau} + \sigma(z_{\tau}) \right)
\]

where $\bar{z}_{\tau}$ is the mean of the four $\tau$ features and $\sigma(z_{\tau})$ is their standard deviation.

$R$ captures the recurrence structure through the primary PC amplitude and the mean and spread of the moment-based timescales.

\section{Participation Ratio Properties}
\label{sec:appendix_pr}

\begin{theorem}
For a Gaussian point cloud $\phi_i \sim \mathcal{N}(\boldsymbol{\mu}, \mathbf{C})$ in $\mathbb{R}^d$, the expected participation ratio is:
\[
\mathbb{E}[\text{PR}] \approx \frac{(\sum_k \lambda_k)^2}{\sum_k \lambda_k^2}
\]
where $\lambda_k$ are the eigenvalues of $\mathbf{C}$.
\end{theorem}

\begin{proof}
For a finite sample of size $N$, the sample covariance matrix $\hat{\mathbf{C}}$ has eigenvalues $\hat{\lambda}_k$ that converge to $\lambda_k$ as $N \to \infty$. The participation ratio is:
\[
\text{PR} = \frac{(\sum_k \hat{\lambda}_k)^2}{\sum_k \hat{\lambda}_k^2}
\]
By the continuous mapping theorem and the convergence $\hat{\lambda}_k \to \lambda_k$, the result follows.
\end{proof}

\section{Adversarial Null Construction}
\label{sec:appendix_adversarial}

The adversarial low-rank null (N6) is designed to produce a manifold with the same participation ratio as the real geometry but with randomized eigenvector orientations.

\begin{enumerate}
\item Compute the target PR from the real $\Phi$-coordinates.
\item Construct eigenvalues $\lambda_k = e^{-k/\alpha}$ where $\alpha$ is chosen so that $\frac{(\sum \lambda_k)^2}{\sum \lambda_k^2} = \text{PR}_{\text{target}}$.
\item Generate random orthogonal eigenvectors $\mathbf{v}_k$ via QR decomposition of Gaussian random matrices.
\item Construct covariance matrix $\mathbf{C} = \sum_k \lambda_k \mathbf{v}_k \mathbf{v}_k^\top$.
\item Sample $\phi_i \sim \mathcal{N}(\mathbf{0}, \mathbf{C})$.
\end{enumerate}

This produces a manifold with matched PR but no meaningful structure in the eigenvector directions.

\section{Causal Destruction Levels}
\label{sec:appendix_destruction}

The causal destruction hierarchy applies progressive transformations to the feature matrix $\mathbf{X} \in \mathbb{R}^{N \times 17}$:

\begin{enumerate}
\item \textbf{Level 0: Original} -- No transformation.
\item \textbf{Level 1: Column shuffle} -- For each feature $j$, randomly permute the values across systems. Destroys feature-system associations.
\item \textbf{Level 2: Row shuffle} -- Randomly permute the rows of $\mathbf{X}$. Destroys system-level feature covariance.
\item \textbf{Level 3: Gaussian copula} -- Transform to Gaussian marginals via rank-based CDF, then back. Preserves marginal distributions but destroys nonlinear structure.
\item \textbf{Level 4: Decorrelated} -- Apply whitening transform to decorrelate features, then remix with random rotation.
\item \textbf{Level 5: White noise} -- Replace $\mathbf{X}$ with i.i.d. Gaussian noise.
\end{enumerate}

\section{Information Geometry}
\label{sec:appendix_info}

The Fisher information matrix for a parametric family $p(\phi; \theta)$ is:
\[
\mathbf{F}(\theta) = \mathbb{E}\left[ \nabla_\theta \log p(\phi; \theta) \nabla_\theta \log p(\phi; \theta)^\top \right]
\]

For the KDE-based density estimate $\hat{\rho}(\phi)$, the Fisher entropy is:
\[
H_F = -\int \log \hat{\rho}(\phi) \, d\phi
\]

The KL divergence from a Gaussian with the same mean and covariance is:
\[
D_{\text{KL}}(\hat{\rho} \| \mathcal{N}) = \int \hat{\rho}(\phi) \log \frac{\hat{\rho}(\phi)}{\mathcal{N}(\phi; \boldsymbol{\mu}, \boldsymbol{\Sigma})} \, d\phi
\]

\end{document}
